% chktex-file 44
% chktex-file 8
\documentclass[10pt,leqno]{amsart}
\usepackage{graphicx}
\baselineskip=16pt

\usepackage{indentfirst,csquotes}

\topmargin= .5cm
\textheight= 20cm
\textwidth= 32cc
\baselineskip=16pt

\evensidemargin= .9cm
\oddsidemargin= .9cm

\usepackage{amssymb,amsthm,amsmath}
\usepackage{xcolor,paralist,hyperref,titlesec,fancyhdr,etoolbox}
\newtheorem{theorem}{Theorem}[]
\newtheorem{definition}[theorem]{Definition}
\newtheorem{example}[theorem]{Example}
\newtheorem{lemma}[theorem]{Lemma}
\newtheorem{proposition}[theorem]{Proposition}
\newtheorem{corollary}[theorem]{Corollary}
\newtheorem{conjecture}[theorem]{Conjecture}

\newcommand{\set}[1]{\left\{#1\right\}}
\newcommand{\abs}[1]{\left|#1\right|}
\newcommand{\paren}[1]{\left(#1\right)}
\newcommand{\Int}[1]{\mathbb Z / #1 \mathbb Z}

% \titleformat{\section}[display]{\normalfont\huge\bfseries\centering}{\centering\chaptertitlename\thechapter}{10pt}{\Large}
% \titlespacing*{\section}{0pt}{0ex}{0ex}

\begin{document}

\title{Twelve lonely runners}
\author[T. Sungkawichai]{Touch Sungkawichai}
\date{\today}
\address{London, United Kingdom}
\email{touchsungkawichai@gmail.com}
\maketitle 

% \let\thefootnote\relax
% \footnotetext{MSC2020: Primary 00A05, Secondary 00A66.}

\begin{abstract}
  The Lonely Runner Conjecture asserts that for any number of runners moving at distinct constant speeds around a unit circle
  from a common starting point, each runner is lonely at some time. A runner is said to be lonely if its distance from every 
  other runner exceeds $1/n$, where $n$ is the total number of runners. 
  Tao\cite{Tao} showed that it suffices to verify the conjecture for finitely many positive integers. 
  Using computational method, Trakultongchai\cite{Trakultongchai} verified the conjecture for up to ten runners.
  In this work, I prove that the conjecture holds for twelve runners, assuming it holds for eleven runners.
\end{abstract}

\bigskip
\section*{Introduction}
The lonely runners is a conjecture stated by Wills~\cite{Wills} that for any $n > 2$, if $n$ runners are 
running with constant speed around a unit circle. Each runner would be eventually lonely. Formally that means for each runner $i$, 
there is a point in time $t$ such that no runner is in distance $\frac{1}{n}$ of runner $i$.
By using the symmetry of the problem, one may fix one runner to be stationary and measure all other speeds relative to it. 
This yields an equivalent formulation as presented in Conjecture~\ref{lonelyconj}.

\begin{conjecture}
  Let $v_1, \dots, v_n$ be distinct positive integers. Then there exists a real number $t$ such that
\[
  \abs{t v_i} \ge \frac{1}{n+1}
\quad \text{for all } i = 1,\dots,n,
\]
  \label{lonelyconj}
\end{conjecture}

The case of a small number of runners, up to seven, has been proven analytically. 
However, as the number of runners increases, the complexity grows exponentially.
In effort to solve this conjecture, Tao\cite{Tao} showed that it is possible to computationally verify the conjecture by only considering integer speeds 
up to $n^{O(n^2)}$, turning a continuous problem into a computational one. 
Later on, Malikiosis et al.\cite{Malikiosis} reduces the search space drastically. They showed that it is enough 
to consider integer speeds up to only $\binom{n+1}{2}^{n-1}$.

Building on the result, Rosenfeld\cite{Rosenfeld} proved the lonely runners conjecture for 8 runners and 
Trakultongchai\cite{Trakultongchai} improved the computational search to assert for 9 and 10 runners.

I improved the computational method used by Rosenfeld and Trakultongchai to conclude that if the conjecture holds for $11$ runners, 
it also holds for 12 runners.

\section*{Rosenfeld's and Trakultongchai's Computational Method}
It was stated by Malikiosis et al.\cite{Malikiosis} that if the $k-1$ lonely runners conjecture holds, 
the $k$ lonely runners can be proven by verifying that there are no 
intregral speed set $V := \set{0, v_1, \dots, v_k}$ that is a counterexample to the conjecture for 
\[ v_1\dots v_k < B_k  \text{where} \quad B_k := \paren{\frac{\binom{k+1}{2}^{k-1}}{k}}^k \]

Rosenfeld introduced the idea of breaking down the bound by small proofs, each of with is a proof that if there is a counterexample, 
then $p \mid v_1 \cdots v_k$. By accumulating enough proofs for a set of primes $P := \set{p_j}$, it follows that 
\[ \prod_{p \in P} p \mid v_1\cdots v_k \] which shows $v_1 \cdots v_k > \prod_{p \in P} p > B_k$

In the work of Trakultongchai\cite{Trakultongchai}, He defined proper and improper solution sets and showed that 
if the set of improper speed sets $I(k, l, p)$ is empty for some value of $l$, then $p \div v_1\cdots v_k$ in any counterexaple.
He argued that such $l$ would need to be a multiple of $k+1$. 

Trakultongchai then proposed an efficient method that allows \textit{lifting} a candidate solution set modulo $p$ 
into a candidate solution set modulo $np$. Allowing computation to use factors of $k+1$ to help speed up the process.
For example, with $k=8$, the idea allows spliting the computation into 3 steps, computing $I(k, 1, p)$ then $I(k, 3, p)$ 
and then $I(k, 9, p)$. 

\section*{Main Work}

The framework introduced in \cite{Trakultongchai} has a bottleneck in the first phase of computation, the 
calculation of $I(k, 1, p)$. 
Notice that for a prime modulus $p$, the objects of interest are speed sets whose elements lie in the multiplicative group
$(\mathbb{Z}/p\mathbb{Z})^\times$.  

Using Trakultongchai's terminology, the sets denoted $I(k,1,p)$ represent the configurations of speeds that could potentially 
violate the Lonely Runner Conjecture modulo $p$.

\begin{definition}
  Let $p$ be a prime number greater than $k+1$.
  A speed set $v := \set{v_1, \dots, v_k} \in ((\Int{p})^\times)^k$ is an improper set modulo $p$, i.e. $v \in I(k, 1, p)$ if for any 
  $t \in \Int{p}$ there exists $i$ such that $\abs{\frac{tv_i}{p}} < \frac{1}{k+1}$. 
\end{definition}

Intuitively, the condition states that for every time parameter $t$ modulo $p$, there is a runner who fails 
to achieve the required separation distance $\frac{1}{k+1}$ from the statonary runner. 
Thus the set is a potential obstruction to the Lonely Runner Conjecture when considered modulo $p$.

\begin{lemma}
For any improper set $\set{v_1, \dots, v_k}$ modulo $p$, there is an improper set $\set{1, u_2, \dots, u_k}$ modulo $p$
and a number $u \in \Int{p}$ such that $\set{u, uu_2, \dots, uu_k} = \set{v_1, \dots, v_k}$.
\end{lemma}

\begin{proof}
  Let $u = v_1$ and $u_i = v_1^{-1} v_i$ for each $i$.  
  Then, $u u_i = v_1 \cdot v_1^{-1} v_i = v_i$. Hence
  \[
      \set{u, u u_2, \dots, u u_k} = \set{v_1, v_2, \dots, v_k}.
  \]
  
  It remains to verify that $\{1,u_2,\dots,u_k\}$ is also an improper set.
  To do that, pick any $t \in (\mathbb{Z}/P\mathbb{Z})^\times$ and set $t' = t v_1^{-1}$.
  Since $\set{v_1, \dots, v_k}$ is an improper set, there exists some index $i$ such that
  \[ \frac{1}{k+1} > \abs{\frac{t' v_i}{p}} = \abs{\frac{t v_1^{-1} v_i}{p}} \]
  If $i = 1$ then $\abs{t/p} < \frac{1}{k+1}$ otherwise $\abs{{t u_i}/{p}} < \frac{1}{k+1}$
  In either case the defining condition holds for the set $\{1,u_2,\dots,u_k\}$, proving that it is also an improper set.
\end{proof}

The lemma shows an equivalent relation of improper sets under multiplication by units. 
Consequently, when enumerating possible counterexamples modulo $p$ it suffices to consider
only normalized sets of the form $\{1,u_2,\dots,u_k\}$.  This reduces the
search space by a factor of $p$, since each equivalence class under scaling
contains $p-1$ representatives.

% TODO: the lift idea 2 -> 4 -> 12 -> 36

In addition to this normalization, I improved the computational approach of
Trakultongchai by refactoring the implementation to support parallel
execution and by introducing more memory-efficient data structures.  The new
implementation remains faithful to the original algorithmic ideas but is
significantly faster in practice.

Using these improvements, I verified the case $k=8$ in approximately
20 seconds and the case $k=9$ in about 11 minutes on a 10-core Apple M4 processor (Macbook Air 2025).

\section*{Result}
Let \[ P_{11} := \set{ \begin{aligned} 
      23, 131, 137, 139, 149, 151, 157, 163, 167, 173, 179, 181, 191, 193, \\ 
      197, 199, 211, 223, 227, 229, 233, 239, 241, 251, 257, 263, 263, 269, \\ 
      271, 277, 281, 283, 293, 307, 311, 313, 317, 331, 337, 347, 349, 353, \\ 
      359, 367, 373, 379, 383, 389, 397, 401, 409, 419, 421, 431, 433, 439, \\ 
      443, 449, 457, 461, 463, 467, 479, 487, 487, 491, 499, 503, 509, 521,\\ 
      523, 541, 547, 557, 563, 569, 571, 577 \\
    \end{aligned}
  } \]

Using the algorithm described in the previous section, the software verifies
that for every prime $p \in P_{11}$ there exists no improper
set modulo $p$. Equivalently, any counterexample to the $12$-runner case of the
Lonely Runner Conjecture must have speeds divisible by each prime in
$P_{11}$. Consequently, if a counterexample exists, its product of speeds must be
divisible by the product
\[ \prod_{p \in P_{11}} p \]

As $B_{11}$, the bound needed to reduce the problem of 12 runners to 11 runners 

\section*{Futher work}
The case of eleven runners remains unsolved. Since 11 is prime, the bottleneck of the computation is not at finding $I(k, 1, p)$ but at 
lifting $I(k, 1, p)$ to $I(k, 11, p)$. Similar arguments can potentially speed up calculation but will require more careful consideration.

\bibliographystyle{plain}
\bibliography{ref}

\end{document}
